\section{Open Questions and Follow-on Probes}

The replication confirmed all closed-form claims of Brown \& Eastin (2018)
at $n=2$ under Pauli noise. Five concrete open questions remain, each with
a bounded next-step experiment on free compute.

\begin{enumerate}

\item \textbf{Scaling to $n=3,4$ qubits.}
\emph{Question.} Do the block-eigenvalue formulas continue to match Stim RB
fits at $n=3,4$, where block-size asymmetry
$(4^n+2^n)/2$ vs $(4^n-2^n)/2$ diverges more strongly and multi-exponential
structure should become more visible?\\
\emph{Basis.} Present work covers only $n=2$; paper's formulas are derived
analytically for arbitrary $n$ with no numerical support at any $n$.\\
\emph{Next step.} Extend \texttt{rb\_replication.py} to $n \in \{3,4\}$
(Stim tableau is tractable for stabilizer RB). Sweep symmetric $p_{\mathrm{dep}}$
and single-axis asymmetric noise; report $\lambda_{\mathrm{fit}}$ vs
theory. Cross-check that walk length 60 remains a sufficient mixing time
by comparing $\lambda_{\mathrm{fit}}$ across walk\_len $\in \{30,60,120,240\}$.

\item \textbf{Hardware-native gate sets (Rydberg, ion-trap M{\o}lmer-S{\o}rensen).}
\emph{Question.} How does the block decomposition play out for the
neutral-atom Rydberg $\{CZ, \text{single-qubit rotations}\}$ set and the
ion-trap $\{\text{MS}_{XX}(\pi/2), \text{single-qubit Cliffords}\}$ set,
which are proper Clifford subgroups distinct from the two the paper covers?\\
\emph{Basis.} Twirl-decomposition proof strategy is generic to any Clifford
subgroup containing the Paulis. MS-XX set may generate a genuinely different
subgroup with different block structure not tabulated in the paper.\\
\emph{Next step.} Brute-force enumerate the Pauli-orbit blocks for the
two native sets at $n=2,3$ (subgroups small enough), then run Stim RB
from $\ket{0\cdots 0}$ and $\ket{+\cdots +}$ and verify the observed
decay matches the enumerated block eigenvalues. Publish a native-set
block-structure table for experimentalists.

\item \textbf{Character-RB layered on the Brown-Eastin blocks.}
\emph{Question.} Can Helsen-style character-RB (using irreducible
characters of the subgroup as measurement weights) resolve all
$\lambda_1,\ldots,\lambda_k$ from a \emph{single} initial state
in one experiment, replacing the paper's one-initial-state-per-block
protocol?\\
\emph{Basis.} Standard RB extracts one $\lambda$ per experiment because
$\rho_0$ overlaps only one block projector; character weighting separates
the blocks. Brown-Eastin block structure is exactly the input a
character-RB needs, but this combination is not discussed in the paper.\\
\emph{Next step.} Compute the irreducible characters of the real-Clifford
and CNOT+Pauli subgroups at $n=2$ (finite, tractable); implement
character-weighted survival in the Stim harness; verify that
$\ket{00}$-only character-RB recovers both $\lambda_1=1$ and
$\lambda_2=0.9472$ under pure-$Z$ noise, and compare sample efficiency.

\item \textbf{ML-boosted decay-fit estimator vs.\ non-linear least squares.}
\emph{Question.} How does an ML-based sequence-to-$\lambda$ regressor
(gradient boosting or a small MLP trained on synthetic Stim runs) perform
against NLSQ in the low-sequence regime $N_{\mathrm{seq}} < 30$ where
the paper's entanglement-infidelity bounds become loose?\\
\emph{Basis.} NLSQ is field-standard but noise-sensitive at short
sequences; recent RB literature (Youssry \emph{et al.}\ 2024) shows
ML fits recover $\lambda$ with 2--3$\times$ lower variance from small
samples by learning the joint $(m, \text{survival})$ distribution
rather than fitting a mean.\\
\emph{Next step.} Generate 10k synthetic runs across the
(subgroup, noise-model, $n$) grid; train xgboost / small MLP with
sklearn on m1; compare RMSE$(\lambda_{\mathrm{fit}}, \lambda_{\mathrm{theory}})$
at $N_{\mathrm{seq}} \in \{10, 30, 60, 100\}$. If ML wins, publish the
trained estimator.

\item \textbf{Coherent (non-Pauli) errors.}
\emph{Question.} Under small gate-independent over-rotation
$R_z(\theta)$ with $\theta \in \{0.01, 0.05, 0.1\}$ rad, do the
block-eigenvalue formulas remain approximate predictors, or does the
non-2-design twirl fail to Pauli-approximate the coherent channel and
introduce block-mixing?\\
\emph{Basis.} Paper explicitly assumes gate-independent Pauli noise;
coherent errors are out of scope. Standard Clifford RB twirls coherent
errors down to $p_{\mathrm{Pauli}} = \sin^2(\theta/2)$ with a correction
$\mathcal{O}(\theta^2)$ that uses the full 2-design property --- the
restricted subgroups lack this, so the closed forms may qualitatively
fail.\\
\emph{Next step.} Extend \texttt{rb\_asym.py} with a coherent-noise
channel (Qiskit statevector for $n=2$; density-matrix tractable);
sweep $\theta \in \{0.001, 0.01, 0.05, 0.1, 0.5\}$ rad; check
whether the decay remains exponential and whether extracted $\lambda$
tracks $p_{\mathrm{Pauli}} = \sin^2(\theta/2)$; quantify any
second-exponential amplitude that indicates block mixing.

\end{enumerate}
